% -*- coding: utf-8 -*-
% Compile with: xelatex pinyin-disambiguation.tex
\documentclass[11pt,a4paper]{article}

\usepackage{fontspec}
\usepackage{xeCJK}
\setCJKmainfont{Noto Serif CJK SC}
\setCJKsansfont{Noto Sans CJK SC}
\setCJKmonofont{Noto Sans Mono CJK SC}

\usepackage{booktabs}
\usepackage{hyperref}
\usepackage{url}
\usepackage[margin=2.5cm]{geometry}
\usepackage{enumitem}
\usepackage{linguex}

\title{Automatic Pinyin Assignment for a Chinese Wordnet\\Using Dictionary Definitions and Gloss Matching}
\author{Francis Bond \and Claude (Anthropic)}
\date{February 2026}

\begin{document}
\maketitle

\begin{abstract}
  We describe a method for automatically assigning Mandarin pinyin
  romanizations to lemmas in the NTU Multilingual Corpus Chinese
  Wordnet.  Our multi-stage pipeline uses CC-CEDICT as the primary
  pronunciation source, resolving ambiguous readings by filtering
  cross-reference-only and weak entries (surname-only, compound-only),
  by matching CEDICT definitions against WordNet English glosses and
  lemmas using word overlap, and by retaining both readings for tone-5
  and same-meaning variant pairs.  Of 234,215 lexical entries, we
  assign pinyin to 230,027 (98.2\%), of which 56,184 are unique
  matches, 477 are disambiguated via gloss and lemma matching, and
  173,134 are assembled character by character.  Only 324
  pronunciation records across 85 lemmas remain genuinely ambiguous.
  We release the remaining ambiguous cases for manual review.
\end{abstract}

\section{Introduction}

The NTU Multilingual Corpus (NTU-MC) Chinese Wordnet contains
234,215 lexical entries linked to Princeton WordNet 3.0 synsets.
While the wordnet records written forms (simplified Chinese
characters), it does not include pronunciation information.  Pinyin
romanization is essential for many downstream applications and for
accessibility.

Adding pinyin is complicated by the fact that many Chinese characters
have multiple readings depending on meaning.  For example, the
character 芥 is read \emph{jiè} when meaning ``mustard'' but
\emph{gài} in the compound 芥蓝 (\emph{gàilán}, ``Chinese
broccoli'').  A simple dictionary lookup returns both readings without
indicating which applies.

We present a multi-stage pipeline that resolves the majority of these
ambiguities automatically, leveraging the structure of both CC-CEDICT
and the wordnet itself.

\section{Resources}

\paragraph{CC-CEDICT} A community-maintained Chinese--English
dictionary containing 120,517 entries with simplified characters,
pinyin, and English definitions.  We use the gzipped release from
MDBG.\footnote{\url{https://www.mdbg.net/chinese/dictionary?page=cedict}}

\paragraph{NTU-MC Chinese Wordnet} The Chinese portion of the NTU
Multilingual Corpus wordnet, containing 234,215 lexical entries mapped
to 116,999 synsets.  Each synset is linked to an English gloss and
English lemmas from Princeton WordNet 3.0 via the NTU-MC database.

\section{Method}

\subsection{Stage 1: Direct CEDICT Lookup}

For each lemma, we look up the written form in CC-CEDICT.  Three
outcomes are possible:

\begin{enumerate}[nosep]
\item \textbf{Unique match}: The word has exactly one reading in
  CEDICT.  This reading is assigned directly.
\item \textbf{Multiple readings}: The word has two or more readings.
  These proceed to Stage~2.
\item \textbf{No match}: The word is not in CEDICT as a whole unit.
  These proceed to character-by-character assembly (Stage~1b).
\end{enumerate}

\subsubsection{Stage 1b: Character-by-Character Assembly}

For words not found in CEDICT, we assemble pronunciation
character by character, using the most frequent reading for each
character as observed across all CEDICT compound entries.  This
frequency-based selection naturally prefers the compound reading of a
character, which is appropriate since these are multi-character words.
The resulting pronunciation is marked with \texttt{notation="composed"}
to indicate lower confidence.

\subsection{Stage 2: Cross-Reference Filtering}

Many CEDICT entries for a character exist only to redirect to a
compound.  For example:

\ex. \textbf{芥} has two CEDICT entries:
  \a. \emph{gài}: ``see 芥蓝 [gài lán]'' (cross-reference only)
  \b. \emph{jiè}: ``mustard'' (real definition)

We identify cross-reference-only readings by checking if all
definitions begin with ``see~\ldots'', ``variant of~\ldots'', or
``old variant of~\ldots''.  After filtering these out, 147 previously
ambiguous words are reduced to a single reading.

\subsection{Stage 2b: Weak-Reading Elimination}

Even among readings with real (non-cross-reference) definitions, some
are unlikely to be the intended reading for a standalone wordnet
lemma.  We classify each reading by its definition types and eliminate
readings that are exclusively:

\begin{itemize}[nosep]
\item \textbf{Surname-only} --- e.g., 曾~\emph{zēng} ``surname
  Zeng'' is eliminated, leaving \emph{céng} ``once; already''.
\item \textbf{Compound-only} (``used in~\ldots'') --- e.g.,
  上~\emph{shǎng} ``used in 上声'' is eliminated, leaving
  \emph{shàng} ``up; upper''.
\item \textbf{Abbreviation-only} --- e.g., 蚌~\emph{bèng} ``abbr.\
  for Bengbu'' is eliminated, leaving \emph{bàng} ``mussel''.
\end{itemize}

\noindent Additionally, readings marked exclusively as colloquial (``(coll.)'')
or Taiwan-specific (``(Tw)'' / ``Taiwan~pr.'') are treated as weak.
This stage resolves an additional ${\sim}90$ words to a
single reading.

\subsection{Stage 3: Gloss and Lemma-Based Disambiguation}

For the remaining words with multiple real readings, we exploit
the fact that each lemma is linked to one or more WordNet synsets,
each of which has an English gloss and a set of English lemmas.  Each
CEDICT reading also comes with English definitions.  We score the
match between CEDICT definitions and the combined gloss and lemma text
using tokenized word overlap.

\paragraph{Scoring} For each reading $r$ with CEDICT definitions
$D_r = \{d_1, d_2, \ldots\}$ and combined synset text $G$ (glosses
$+$ English lemmas):

\begin{enumerate}[nosep]
\item Tokenize all texts: lowercase, remove punctuation and stop words.
\item Apply rough suffix stripping (\emph{-tion}, \emph{-ing},
  \emph{-ed}, \emph{-ly}, etc.) to improve recall.
\item Compute $\mathrm{score}(r) = |T(D_r) \cap T(G)|$, where $T(\cdot)$
  denotes the token set.
\end{enumerate}

If one reading scores strictly higher than all others, it is selected.
Including English lemmas (e.g., \emph{play}, \emph{strum},
\emph{consequence}) alongside glosses substantially improves recall,
since lemmas often share vocabulary with CEDICT definitions more
directly than the longer gloss text.

\paragraph{Examples}

\begin{table}[ht]
\centering
\begin{tabular}{llll}
  \toprule
  Word & Readings & Synset lemmas + gloss & Selected \\
  \midrule
  弹 & \emph{tán} `to \textbf{play}; to \textbf{flick}' & lemmas: \emph{\textbf{play}, \textbf{flip}, \textbf{flick}} & \emph{tán} \\
     & \emph{dàn} `bullet; shell' & gloss: ``\textbf{play} on an instrument'' & \\[4pt]
  打 & \emph{dǎ} `to \textbf{hit}; to \textbf{strike}' & lemmas: \emph{\textbf{hit}, \textbf{strike}} & \emph{dǎ} \\
     & \emph{dá} `dozen (loanword)' & gloss: ``\textbf{hit} with a missile \ldots'' & \\[4pt]
  结果 & \emph{jiéguǒ} `\textbf{outcome}; \textbf{result}' & lemmas: \emph{\textbf{result}, \textbf{consequence}} & \emph{jiéguǒ} \\
       & \emph{jiēguǒ} `to bear \textbf{fruit}' & gloss: ``a phenomenon that follows \ldots'' & \\[4pt]
  行 & \emph{háng} `\textbf{row}; \textbf{line}' & lemmas: \emph{\textbf{row}, \textbf{column}} & \emph{háng} \\
     & \emph{xíng} `to walk; to travel' & gloss: ``an arrangement of objects \ldots'' & \\
  \bottomrule
\end{tabular}
\caption{Examples of gloss and lemma-based disambiguation.  The
  reading whose CEDICT definitions share content words (bold) with
  the synset lemmas and gloss is selected.}
\label{tab:examples}
\end{table}

\subsection{Stage 4: Tone-5 Variant Pairs}

When two readings of a word differ only in whether a syllable carries
a neutral tone (tone~5) versus a full tone, both readings are retained
as valid pronunciations.  For example, 东西 has readings
\emph{dōngxi} (``thing; stuff'') and \emph{dōngxī} (``east and
west''); both are kept.  This stage resolves ${\sim}48$ previously
ambiguous entries.

\subsection{Stage 5: Same-Meaning Variants}

Some words have multiple readings that share overlapping definitions
in CEDICT --- these are genuine pronunciation variants of the same
meaning.  For example, 尽量 can be read either \emph{jǐnliàng} or
\emph{jìnliàng}, both meaning ``as much as possible''; 削 can be
read \emph{xiāo} or \emph{xuē}, both meaning ``to cut; to pare''.
When we detect that two readings share at least one common definition
(case-insensitive), we retain both.  This resolves ${\sim}11$
entries.

\subsection{Stage 6: Taiwan Cross-References}

When one of two remaining readings is annotated in CEDICT with
``Taiwan~pr.\ [\ldots]'' pointing to the other reading, we keep both
pronunciations: the mainland reading with no notation and the Taiwan
reading with \texttt{notation="Taiwan"}.  For example, 血 has
\emph{xuè} (standard) and \emph{xiě} (Taiwan~pr.), both meaning
``blood''; both are retained.  This stage resolves 9~entries.

\subsection{Pronunciation Notes}

The \texttt{notation} attribute of the WN-LMF \texttt{Pronunciation}
element encodes the provenance of each reading:

\begin{itemize}[nosep]
\item No \texttt{notation} --- resolved (unique CEDICT match or
  disambiguated).
\item \texttt{notation="composed"} --- character-by-character assembly.
\item \texttt{notation="unresolved"} --- multiple readings retained.
\item \texttt{notation="Taiwan"}, etc.\ --- regional or register
  variant from a CEDICT pronunciation note (e.g., ``Taiwan~pr.'',
  ``colloquial~pr.'').
\end{itemize}

\section{Results}

Table~\ref{tab:results} summarizes the outcome of the pipeline across
all 234,215 lexical entries.

\begin{table}[ht]
\centering
\begin{tabular}{lrr}
  \toprule
  Category & Entries & \% \\
  \midrule
  No pinyin (no CEDICT data) & 4,188 & 1.8 \\
  \midrule
  Unique CEDICT match & 56,184 & 24.0 \\
  \quad incl.\ cross-reference filtered & 147 & --- \\
  \quad incl.\ weak-reading eliminated & ${\sim}90$ & --- \\
  Gloss/lemma-disambiguated & 477 & 0.2 \\
  Tone-5 / same-meaning (both kept) & 72 & --- \\
  Taiwan cross-references (both kept) & 9 & --- \\
  Character-by-character (composed) & 173,134 & 73.9 \\
  \midrule
  \emph{Total with pinyin} & \emph{230,027} & \emph{98.2} \\
  \midrule
  Remaining ambiguous (unresolved) & 324 & --- \\
  \quad (85 unique lemmas, 151 entries) & & \\
  \bottomrule
\end{tabular}
\caption{Pinyin assignment results.  Remaining ambiguous entries have
  multiple pronunciation records attached.  Tone-5 pairs and
  same-meaning variants (${\sim}59$ entries) are counted under the
  unique match total since both readings are retained.}
\label{tab:results}
\end{table}

\paragraph{Disambiguation breakdown}  Of the words with multiple
CEDICT readings, 147 were resolved by cross-reference filtering
(Stage~2), ${\sim}90$ by weak-reading elimination (Stage~2b), 477 by
gloss and lemma matching (Stage~3), ${\sim}48$ by tone-5 pair
retention (Stage~4), ${\sim}11$ by same-meaning detection
(Stage~5), and 9 by Taiwan cross-reference retention (Stage~6),
leaving 85 unique words (324 pronunciation records across 151 entries)
unresolved.

The unresolved cases fall into several categories:

\begin{itemize}[nosep]
\item \textbf{Function words} (e.g., 不是, 了, 么): Definitions are
  too abstract for word overlap to distinguish readings.
\item \textbf{Genuinely sense-dependent} (e.g., 中 `middle' vs.\
  `to hit'; 乐 `happy' vs.\ `music'): The correct reading depends on
  which synset the entry is linked to, but neither CEDICT definition
  overlaps the English gloss or lemmas.
\end{itemize}

\subsection{Additional Processing: Particle Lemmas}

The NTU-MC database includes 29,830 entries where the Chinese lemma
contains \emph{+的} or \emph{+地}, indicating adjectival or adverbial
derivation (e.g., 高兴+的 ``happy [adj]'').  We normalize these by
removing the \emph{+} sign and marking all senses of the entry as
\texttt{lexicalized="false"} in the WN-LMF output, following the
convention that these are productive morphological forms rather than
independent lexical items.

\section{Discussion}

The combined pipeline reduces ambiguous lemmas from over 1,000 to
85.  Gloss and lemma matching resolves 477 entries; the structural
heuristics (cross-reference filtering, weak-reading elimination,
tone-5 pairs, same-meaning variants, and Taiwan cross-references)
resolve a further ${\sim}240$.
Including English lemmas alongside glosses in the scoring was a
significant improvement, nearly halving the number of unresolved
cases.

While more sophisticated methods such as word embeddings or
cross-lingual models could potentially improve this, the marginal
gains are likely small: most remaining cases involve function words or
genuinely polyphonic characters where English glosses and lemmas
provide little discriminative signal.

Regional pronunciation variants from CEDICT (e.g., ``Taiwan~pr.'')
are preserved as \texttt{notation} annotations, allowing downstream
applications to present or filter variant readings as needed.

The 85 remaining ambiguous lemmas are released as a tab-separated
file for manual verification.  Each row contains the word, reading,
overlap score, CEDICT definitions, and WordNet glosses, facilitating
efficient human review.

\section{Conclusion}

We have presented a lightweight, resource-efficient method for
assigning pinyin to a Chinese wordnet that requires no training data
or neural models---only a bilingual dictionary (CC-CEDICT) and the
English glosses and lemmas already present in the wordnet.  The method
achieves 98.2\% coverage with high precision on the resolved cases,
and produces a structured list of remaining ambiguities for human
review.

\appendix
\section{Random Examples}

This appendix shows 10 randomly selected entries from each of the
three main output categories.

\subsection{Resolved (unique or disambiguated)}

These entries received a single pinyin reading, either from a unique
CEDICT lookup or via gloss/lemma disambiguation.

\begin{enumerate}[nosep]
\item 杂乱无章 \emph{zá luàn wú zhāng} `disorderly'
\item 幼儿 \emph{yòu'ér} `infant'
\item 克娄巴特拉 \emph{kè lóu bā tè lā} `Cleopatra'
\item 倒数 \emph{dào shù} `countdown'
\item 俄克拉荷马 \emph{é kè lā hé mǎ} `Oklahoma'
\item 黑天 \emph{hēi tiān} `Krishna'
\item 说道 \emph{shuō dào} `to say'
\item 闽南语 \emph{mǐn nán yǔ} `Hokkien'
\item 军国主义 \emph{jūn guó zhǔ yì} `militarism'
\item 反语 \emph{fǎn yǔ} `irony'
\end{enumerate}

\subsection{Character-by-character (composed)}

These entries were not found in CEDICT as whole words; their
pronunciation was assembled from individual character frequencies.

\begin{enumerate}[nosep]
\item 复燃室 \emph{fù rán shì} `afterburner chamber'
\item 铀矿 \emph{yóu kuàng} `uranium ore'
\item 地方教育委员会 \emph{dì fāng jiào yù wěi yuán huì} `local education committee'
\item 连水陆路 \emph{lián shuǐ lù lù} `combined water-land route'
\item 劳役马 \emph{láo yì mǎ} `draft horse'
\item 马来麝 \emph{mǎ lái shè} `Malay musk deer'
\item 令人安慰的 \emph{lìng rén ān wèi de} `reassuring'
\item 燕子水仙 \emph{yàn zi shuǐ xiān} `jonquil'
\item 倾角罗盘 \emph{qīng jiǎo luó pán} `dip compass'
\item 铀矿 \emph{yóu kuàng} `uranium mine'
\end{enumerate}

\subsection{Unresolved (multiple readings)}

These entries retain two or more candidate pronunciations, marked with
\texttt{notation="unresolved"}.

\begin{enumerate}[nosep]
\item 通 \emph{tōng} / \emph{tòng}
\item 教学 \emph{jiāo xué} / \emph{jiào xué}
\item 宿 \emph{xiǔ} / \emph{xiù}
\item 好玩 \emph{hǎo wán} / \emph{hào wán}
\item 台 \emph{tāi} / \emph{tái}
\item 斗 \emph{dǒu} / \emph{dòu}
\item 勾 \emph{gōu} / \emph{gòu}
\item 渐 \emph{jiān} / \emph{jiàn}
\item 追 \emph{duī} / \emph{zhuī}
\item 将 \emph{jiāng} / \emph{jiàng} / \emph{qiāng}
\end{enumerate}

\end{document}
